\chapter{Conclusion}
\label{ch:conclusion}

\section{Summary of findings}
\label{sec:concl-summary}
Camera traps and laboratory rigs now record wildlife video far faster than anyone can annotate it, and that
annotation step is what limits how much ecology and behavioural science actually gets done. Promptable
trackers offer a way past it. Name an animal in ordinary words, and the model finds and follows every one that
appears, without ever having been trained on that species. The difficulty is that the ability is uneven, and
uneven in ways nobody can currently anticipate. A team pointing such a model at a new species in a new place
has no principled basis for deciding whether to trust what comes back, which leaves two unattractive options:
trust it blindly, or fall back on the manual annotation the model was meant to replace.

This dissertation asked whether that judgement can be made in advance. Can a promptable tracker's zero-shot
accuracy on an unseen species and place be predicted before the model is run, from properties that require no
labelling of the target? It asked the question separately for the two halves of the task --- finding the
animal, and following it --- because the two can fail for different reasons.

The answer is no, and it survives every criterion applied to it. No distance's coefficient interval excludes
zero. None beats a mean predictor out of sample, on either a species or a location hold-out. Together the four
explain about $3\%$ of the variation in detection. The null holds when the encoder, the crop, the prompt and
the distance functional are each swapped in turn. It holds for a second family of label-free properties,
intrinsic scene difficulty, whose apparent effect collapses into animal size. It holds when familiarity is
read from the tracker's own feature space instead of from an external reference. It holds on a second,
independent dataset, and on two architecturally distinct trackers substituted for SAM\,3. And it is not a
failure of power: the same procedure detects animal size out of sample wherever size is present.

Only that one property predicts detection, and only modestly --- the nuisance the analysis introduced in order
to control for it. Association carries almost no variance distinct from detection on this data, so the
decomposition that motivated the design turned out to be a detection story. What remains is a rigorous,
decomposed, cross-dataset and cross-model null.

\section{Contributions}
\label{sec:concl-contributions}
To our knowledge this is the first study to test, out of sample, whether a promptable video tracker's
zero-shot transfer is predictable from label-free signals known before the model is run, and the first to
decompose that question into detection and association.

What the field gains is threefold. The formulation and the protocol are reusable: any future before-running
estimator can be held to the same whole-group hold-out bar that these well-motivated distances failed, which
is what makes a negative answer useful rather than merely discouraging. The explanation is specific rather
than vague: the apparent signals reduce to an animal-size confound, so any future study measuring visual
novelty across species must control for apparent size or it will rediscover the same artefact. And the caution
is practical: proximity to familiar training data is not a safety guarantee, so reliability has to be checked
rather than assumed.

\section{Future work}
\label{sec:concl-future}
Three directions remain open.

The first is a fuller representational probe. The familiarity proxy tests only whether SAM\,3's features
separate the species, and finds that they do by re-encoding the visual distance
(Section~\ref{sec:familiarity}). What those features actually encode --- whether taxonomy is represented
structurally --- is untested, as is a proxy anchored on the coverage of the true, undisclosed pretraining
corpus.

The second is a broader model swap: more trackers, and ideally a larger fair dataset that is no tracker's home
turf. SA-FARI's larger cell count could not be reused here, for the reason given in
Section~\ref{sec:model_swap}.

The third is targets with richer association structure. In crowded, multi-object footage the
detection--association decomposition would have considerably more to distinguish than it did here.
